\section{Introduction}
\label{sec:introduction}

Suppose an empirical researcher uses a news article to form a signal
before the market opens and measures the stock's return from the open to
the close. A second researcher can observe the same price change and
apply the same estimator, but reconstructing the signal requires the
version of the article that was available at the time and an identifier
for the security to which the text refers. Because an article can be
updated after its initial publication, even a record of the publisher's
first-publication time leaves open which text the original analysis used.
The data behind the signal include these decisions as well as the words
on the page.

Access to proprietary news is one obstacle to recovering those decisions.
\citet{heston2017}, for example, use a large Reuters corpus and a
proprietary sentiment measure to study how stock prices respond to news.
The data and measurements make the question tractable. A researcher who
has the estimating equation but lacks access to the article versions or
their sentiment annotations must still recover the inputs before
repeating the analysis or changing the way sentiment is measured. A
publicly accessible web archive can supply an observed version of a
story, and an archived copy gives the research team a common input even
when the publisher changes or removes the live page.

We ask when those copies can support a reconstructible sample for
financial research, and how the rules that connect captures to issuers
and trading times change the resulting measurements. We construct
CCFinance from a frozen selection of Common Crawl News, preserve
available source and timing evidence, link extracted company mentions to
dated issuer records, and evaluate the resulting populations under
explicit eligibility rules. A link records which company the text names.
A separate decision about company focus determines whether the same text
is suitable for studying that firm's news. Recording the decisions makes
it possible to compare what a researcher gains and loses by tightening
the definition of an observation.

Three exercises test successive requirements on those observations. An
equity-market volatility tracker requires a denominator of eligible
articles for each publisher and month. A co-mention network requires
articles linked to pairs of banks. An overnight return test also
requires evidence that the selected article was available before the
trading interval and enough securities with qualifying signals to form
portfolios. We first compare the types of source evidence available in
prior corpora, then measure the support produced by these requirements.
